\chapter{Stochastic Volatility}
\label{chap:stochastic_volatility}
Volatility is a fundamental concept in financial econometrics and is commonly used as a measure of the variability or uncertainty of asset returns. In its most basic form, volatility refers to the standard deviation of returns over a given time horizon, and therefore quantifies the magnitude of fluctuations around the expected return.

Many models assume that volatility is constant. In practice, however, financial markets exhibit periods of high volatility, for example, during financial crises, and periods of low volatility, usually during stable economic conditions. Stochastic volatility models capture this behavior by allowing volatility to vary over time according to a stochastic process.

This chapter begins by introducing volatility and then develops the state-space formulation used throughout the thesis, with particular emphasis on the Markov switching stochastic volatility (MSSV) model.

\section{Volatility}
\label{sec:volatility}
Let $P_t$ denote the price of an asset (stock, index, commodity, etc.) at time $t$.

\begin{definition}
    The \emph{log return} of the asset between two consecutive time points is defined as
    \begin{equation}
        Y_t = \log(P_t) - \log(P_{t-1}) = \log\left(\frac{P_t}{P_{t-1}}\right).
    \end{equation}
\end{definition}

The advantage of using log returns rather than simple returns lies in their additive structure over time and their natural treatment of compounding effects \cite{fan2017elements}. Throughout this thesis, the term return will also refer to a log return.

Figure \ref{fig:log_returns} shows the log returns of the S\&P 500 index from 1980 till mid 2026. Clearly the data exhibits heteroscedasticity, as the variance is not constant over time. In fact, we can see clusters of high and low variability, which will motivate the use of multiple regimes to model the volatility process.
\begin{figure}[h]
    \centering
    \includegraphics[width=0.99\textwidth]{figures/log_returns_sp500.png}
    \caption{Log returns of the S\&P 500 index from 1980 till mid 2026.}
    \label{fig:log_returns}
\end{figure}    

While returns are directly observable, volatility is a characteristic of the underlying return distribution and is not directly observable from a single realization. At each time point $t$, only one return is observed, whereas volatility is defined through the variability of the entire return-generating distribution. Consequently, time-varying volatility must be inferred from data rather than measured directly.

\begin{definition}
    Let $\mathcal{F}_{t-1}$ denote the information set containing all relevant information available up to time $t-1$. The \emph{conditional mean} and \emph{conditional variance} of the return are defined as
    \begin{align}
        \mu_t &= E[Y_t \mid \mathcal{F}_{t-1}], \\
        \sigma_t^2 &= \operatorname{Var}(Y_t \mid \mathcal{F}_{t-1}) = E[(Y_t - \mu_t)^2 \mid \mathcal{F}_{t-1}],
    \end{align}
    and $\sigma_t$ is referred to as the \emph{conditional volatility} \cite{andersen2006volatility}.
\end{definition}

In many financial applications, the conditional mean is assumed to be negligible relative to the variability of returns, so that the focus is primarily on modelling the conditional volatility process.

So far, we have not specified the time scale of the returns. For example, $Y_t$ could represent daily, weekly, or monthly returns, and $t$ could therefore index days, weeks, or months.
In this thesis, we focus on daily returns, specifically close-to-close daily returns, meaning that $P_t$ is the closing price of the asset on day $t$\footnote{One could also consider open-to-open returns, where $P_t$ is the opening price of the asset on day $t$.}.
The goal is then to estimate the daily conditional volatilities $\sigma_t$. 

\subsection{Realized variance}
\label{sec:realized_variance}
When high-frequency intraday data are available, the latent conditional variance can be approximated using realized variance. Consider a day $t$, divided into $m$ equally spaced intraday intervals. Let $Y_{t,i}^{(m)}$ denote the return over the $i$-th intraday interval, $i = 1, \ldots, m$, so that the corresponding daily return can be written as
\begin{equation}
    Y_t = \sum_{i=1}^{m} Y_{t,i}^{(m)}.
\end{equation}

\begin{definition}
    The \emph{realized variance} of day $t$ at sampling frequency $m$ is defined as
    \begin{equation}
        RV_t^{(m)} = \sum_{i=1}^{m} \left(Y_{t,i}^{(m)}\right)^2.
    \end{equation}
\end{definition}

It is often reasonable to assume that 
\begin{align}
    E\left[Y_{t,i}^{(m)} \mid \mathcal{F}_{t-1}\right] &\approx 0, \\
    \operatorname{Cov}\left(Y_{t,i}^{(m)}, Y_{t,j}^{(m)} \mid \mathcal{F}_{t-1}\right) &\approx 0, \qquad \text{for} \; i \neq j.
\end{align}
Under these assumptions, we have
\begin{align}
    \sigma_t^2 = \operatorname{Var}(Y_t \mid \mathcal{F}_{t-1}) &= \operatorname{Var}\left(\sum_{i=1}^{m} Y_{t,i}^{(m)} \,\middle|\, \mathcal{F}_{t-1}\right) \\
    &\approx \sum_{i=1}^{m} \operatorname{Var}\left(Y_{t,i}^{(m)} \mid \mathcal{F}_{t-1}\right) \\
    &\approx \sum_{i=1}^{m} E\left[\left(Y_{t,i}^{(m)}\right)^2 \,\middle|\, \mathcal{F}_{t-1}\right] \\
    &= E\left[RV_t^{(m)} \mid \mathcal{F}_{t-1}\right].
\end{align}
Thus, under these assumptions, realized variance provides an approximately unbiased estimator of the conditional variance. 
Moreover, as the sampling frequency increases, the realized variance generally becomes a more precise measure of the underlying volatility, since the estimation error decreases with $m$.
In practice, however, $m$ must not be too large, to prevent intraday returns from becoming correlated due to market microstructure effects \cite{hansen2005forecast}.

\subsection{Overnight returns}
\label{sec:overnight_returns}
Many assets are not traded continuously, as markets are closed overnight, over weekends, and on holidays. We therefore decompose the daily return into overnight (close-to-open) and intraday (open-to-close) components,
\begin{equation}
    Y_t = \log\left(\frac{P_t^{close}}{P_{t-1}^{close}}\right) = \log\left(\frac{P_t^{open}}{P_{t-1}^{close}}\right) + \log\left(\frac{P_t^{close}}{P_t^{open}}\right) = Y_t^{CO} + Y_t^{OC}
\end{equation}
where $Y_t^{CO}$ is the overnight (close-to-open) return and $Y_t^{OC}$ is the intraday (open-to-close) return. Here we see the usefulness of the additive structure of log returns.
From this it naturally follows that the daily variance can be decomposed into an overnight and an intraday component, if we assume that the two components are uncorrelated:
\begin{align}
    \sigma_t^2 = \operatorname{Var}(Y_t \mid \mathcal{F}_{t-1}) &= \operatorname{Var}(Y_t^{CO} + Y_t^{OC} \mid \mathcal{F}_{t-1}) \\
    &= \operatorname{Var}(Y_t^{CO} \mid \mathcal{F}_{t-1}) + \operatorname{Var}(Y_t^{OC} \mid \mathcal{F}_{t-1}).
\end{align}

While the intraday variance can be estimated from high-frequency data using realized variance $RV_t^{(m)}$, the overnight variance is difficult to estimate reliably from a single observation\footnote{In case when high-frequency overnight trading data is available, the overnight variance can be approximated in the same way using intranight squared returns, just like the intraday variance \cite{martens2002measuring}.}.
One solution is to simply add the squared overnight return to $RV_t^{(m)}$. The variance estimator is then given by
\begin{equation}
    \hat{\sigma}_t^2 = RV_t^{(m)} + (Y_t^{CO})^2.
\end{equation}
This approach is problematic, as the squared overnight return is a very noisy estimator of the overnight variance \cite{martens2002measuring}. 

A better approach is to scale the intraday realized variance, which gives the estimator
\begin{equation}
    \hat{\sigma}_t^2 = c RV_t^{(m)}, \qquad c = \frac{\frac{1}{T}\sum_{t=1}^T \left(Y_t - \hat{\mu}_t\right)^2}{\frac{1}{T}\sum_{t=1}^T RV_t^{(m)}},
\end{equation}
where $\hat{\mu}_t$ is the sample mean of the daily returns and $T$ is the total number of days in the sample.
Here, $c$ measures the historical ratio of total close-to-close variance to intraday variance. Equivalently, if $\kappa$ denotes the historical ratio of overnight to intraday variance, then $c=1+\kappa$.
Thus, the missing overnight variance is estimated as a proportion $\kappa$ of the intraday realized variance, rather than from the noisy overnight return of each individual day. Such an approach is valid under the assumption that intraday and overnight returns are uncorrelated.
Weekends and holidays are naturally treated as extended overnight periods \cite{hansen2005forecast,martens2002measuring}.

Overall, when high-frequency intraday observations are available, they provide a direct and increasingly precise proxy for the underlying daily variance. In empirical applications, realized variance can therefore be used in place of the unobserved conditional variance, which provides a way to evaluate the performance of volatility models. 

\section{State-space models}
\label{sec:state_space_models}
A state-space model describes a dynamical system in which an unobserved (latent) stochastic process drives the evolution of an observed process. The central idea is that the system evolves through hidden states, while only noisy measurements of those states are available.

Let $\{X_t\}_{t \geq 0}$ denote the sequence of random variables representing the latent state process, let $\{Y_t\}_{t \geq 1}$ denote the sequence of random variables representing the observed process, and let $T$ denote the time horizon, hence $t \in \{0, 1, \ldots, T\}$\footnote{The latent states are indexed from $0$ to $T$, whereas the observed returns are indexed from $1$ to $T$. This emphasizes that the initial state $X_0$ is the prior and is not associated with any observation. The indexing could be defined differently (e.g., with observations also starting at $t=0$), but we adopt this convention throughout.}. 
$X_t$ takes values in a measurable state space $(\mathcal{X}, \mathcal{B}(\mathcal{X}))$ and $Y_t$ takes values in $(\mathcal{Y}, \mathcal{B}(\mathcal{Y}))$. Together, $\{(X_t, Y_t)\}_{t \geq 1}$ forms a stochastic process on the joint space $(\mathcal{X} \times \mathcal{Y}, \mathcal{B}(\mathcal{X}) \otimes \mathcal{B}(\mathcal{Y}))$ \cite{chopin2020smc}.

\begin{definition}
    \label{def:state_space_model}
    A \emph{state-space model} with parameters $\theta \in \Theta$ is a pair of processes $\{X_t\}_{t \geq 0}$ and $\{Y_t\}_{t \geq 1}$ such that
    \begin{enumerate}
        \item the latent process $\{X_t\}_{t \geq 0}$ is a Markov process,
        \item the observations $\{Y_t\}_{t \geq 1}$ are conditionally independent given the latent process, and $Y_t$ depends on the latent process only through $X_t$,
    \end{enumerate}
    specified by the densities
    \begin{align}
        X_0 &\sim \chi(x_0 \mid \theta), \\
        X_t \mid X_{t-1}=x_{t-1} &\sim p(x_t \mid x_{t-1}, \theta), \\
        Y_t \mid X_t=x_t &\sim f(y_t \mid x_t, \theta),
    \end{align}
    where $\chi$ is the \emph{initial density}, $p$ is the \emph{transition density}, and $f$ is the \emph{observation} or \emph{emission density} \cite{chopin2020smc, doucet2001sequential,tingstrom2015smc2sv}.
\end{definition}

\begin{remark}
    The conditional distribution of $X_t \mid X_{t-1}$ need not admit a density in general. A proper way would be to define the state-space model in a more measure focused way using probability kernels \cite{chopin2020smc}.
    However, we will not deal with this more general definition in this thesis, as our models admit densities.
\end{remark}

Sometimes, the state-space model is defined in even more general terms, where the latent process and observation process are also dependent on the control input $U_t$, which constitutes another sequence of random variables \cite{murphy2012machine}. Such a definition is more common in control theory and engineering. The control inputs can be thought of as external factors that influence the system dynamics and observations. 
In this thesis, we will not consider control inputs, as we are primarily interested in modelling the latent volatility process and its relationship with observed returns.

The main goal of state-space models is to estimate the latent states based on the observed data. With $X_{i:j} = \{X_i, X_{i+1}, \ldots, X_j\}$ we denote the sequence of random variables from time $i$ to $j$, and with $x_{i:j} = \{x_i, x_{i+1}, \ldots, x_j\}$ we denote their realizations.

\begin{definition}
    \label{def:distributions_of_interest}
    Let $p(x_{0:t}, \theta \mid y_{1:t})$ be the joint posterior distribution of the latent states and model parameters given the observed data. We define
    \begin{itemize}
        \item the \emph{smoothing distribution}
        \begin{equation}
            \label{eq:smoothing_distribution}
            p(x_{0:t} \mid y_{1:t}) = \int p(x_{0:t}, \theta \mid y_{1:t})\,\mathrm{d}\theta,
        \end{equation}
        \item the \emph{filtering distribution}
        \begin{equation}
            \label{eq:filtering_distribution}
            p(x_t \mid y_{1:t}) = \int p(x_{0:t}, \theta \mid y_{1:t})\,\mathrm{d}x_{0:t-1}\,\mathrm{d}\theta,
        \end{equation}
        \item the \emph{predictive distribution}
        \begin{equation}
            \label{eq:predictive_distribution}
            p(x_{t+1} \mid y_{1:t}) = \int p(x_{t+1} \mid x_t, \theta) p(x_{0:t}, \theta \mid y_{1:t})\,\mathrm{d}x_{0:t}\,\mathrm{d}\theta,
        \end{equation}
        \item the \emph{posterior distribution of the parameters}
        \begin{equation}
            \label{eq:posterior_distribution_parameters}
            p(\theta \mid y_{1:t}) = \int p(x_{0:t}, \theta \mid y_{1:t})\,\mathrm{d}x_{0:t}.
        \end{equation}
    \end{itemize}
\end{definition}

When the parameters $\theta$ are known, the smoothing, filtering and predictive distributions are defined analogously, by conditioning on $\theta$ instead of integrating it out, i.e., $p(x_{0:t} \mid y_{1:t}, \theta)$, $p(x_t \mid y_{1:t}, \theta)$ and $p(x_{t+1} \mid y_{1:t}, \theta)$.

\begin{definition}
    A \emph{hidden Markov model} (HMM) is a state-space model whose latent state takes values in a finite set $\mathcal{X} = \{1, 2, \ldots, K\}$, where $K$ is the number of states.
\end{definition}

Although there is some discrepancy in the literature, and the term hidden Markov model is sometimes used for state-space models in general, we adopt the definition above \cite{chopin2020smc}.

% \subsection{Linear Gaussian state space model}
% \label{subsec:linear_gaussian_ssm}
% An important class of state space models is the linear Gaussian state space model (LG-SSM), in which both the transition and observation equations are linear and Gaussian. Mathematically, the model is defined as
% \begin{align}
%     X_t &= A_t X_{t-1} + \epsilon_t, \quad \epsilon_t \sim \mathcal{N}(0, Q_t), \\
%     Y_t &= C_t X_t + \delta_t, \quad \delta_t \sim \mathcal{N}(0, R_t),
% \end{align}
% where $A_t, C_t$ are matrices that define the dynamics, and $Q_t$ and $R_t$ are covariance matrices of the process and observation noise, respectively. Note that if matrices $A_t, C_t$ are time independent, the model is said to be stationary. 
% In literature we often find the LG-SSM defined in a more general form with control inputs, by adding $B_t U_t$ to the transition equation and $D_t U_t$ to the observation equation, where $U_t$ is a control input \cite{murphy2012machine}. However, as mentioned in Section \ref{sec:state_space_models}, we will not consider control inputs in this thesis.

% We can write the model in a more density focused notation:
% \begin{align}
%     X_t | X_{t-1}=x_{t-1} &\sim \mathcal{N}(A_t x_{t-1}, Q_t), \\
%     Y_t | X_t=x_t &\sim \mathcal{N}(C_t x_t, R_t).
% \end{align}
% Let $\theta$ denote the collection of model parameters.
% If at time $t=0$, the initial state is Gaussian as well:
% \begin{equation}
%     X_0 \sim \mathcal{N}(\mu_0, Q_0),
% \end{equation}
% then all subsequent filtering distributions $p(x_t \mid y_{1:t}, \theta)$ are also Gaussian.
% Using the algorithm known as the Kalman filter, we can compute the parameters of these Gaussian filtering distributions in a closed form \cite{murphy2012machine}. 

\section[Markov switching stochastic volatility model]{Markov switching stochastic volatility\\ model}
\label{sec:mssv}
The Markov switching stochastic volatility (MSSV) model is the main model of interest in this thesis. 
It is a state-space model in which the hidden state is the tuple $X_t = (H_t, S_t)$. The log volatility $H_t$ is continuous, while the regime indicator $S_t$ is discrete and takes values in $\{1, 2, \ldots, K\}$, where $K$ is the number of regimes \cite{so1998stochastic}. 

\begin{definition}
    \label{def:mssv}
    The \emph{Markov switching stochastic volatility model} with $K$ regimes is the state-space model with latent state $X_t = (H_t, S_t)$ taking values in
    \begin{equation}
        \mathcal{X} = \mathbb{R} \times \{1, 2, \ldots, K\},
    \end{equation}
    equipped with the product $\sigma$-algebra $\mathcal{B}(\mathbb{R}) \otimes 2^{\{1, 2, \ldots, K\}}$, and the following distributions \cite{so1998stochastic, carvalho2007simulation}.
    \begin{itemize}
        \item \textbf{Observation distribution:}
        \begin{equation}
            Y_t \mid X_t=(h_t, s_t) \sim \mathcal{N}\left(0, \exp\left(h_t\right)\right).
        \end{equation}
        \item \textbf{Transition distribution:}
        \begin{align}
            H_t \mid H_{t-1}=h_{t-1}, S_t=s_t &\sim \mathcal{N}(\mu_{s_t} + \phi(h_{t-1} - \mu_{s_t}), \eta^2), \\
            S_t \mid S_{t-1}=s_{t-1} &\sim \text{Categorical}([p_{s_{t-1},1}, \ldots, p_{s_{t-1},K}]),
        \end{align}
        where the regime $S_t$ is conditionally independent of $H_{0:t-1}$ and $S_{0:t-2}$ given $S_{t-1}$, and the log volatility $H_t$ is conditionally independent of $H_{0:t-2}$ and $S_{0:t-1}$ given $H_{t-1}$ and $S_t$.
        \item \textbf{Initial distribution:}
        \begin{align}
            H_0 \mid S_0=s_0 &\sim \mathcal{N}\left(\mu_{s_0}, \frac{\eta^2}{1 - \phi^2}\right), \\
            S_0 &\sim \text{Categorical}\left(\frac{1}{K}, \frac{1}{K}, \ldots, \frac{1}{K} \right).
        \end{align}
    \end{itemize}
    Here $\mathcal{N}$ denotes the normal distribution with the given mean and variance, $\mu_{s}$ is the regime-specific mean of regime $s$, $\phi$ is the persistence parameter, $\eta^2$ is the variance of the log-volatility innovation process, and $p_{ij} = \mathbb{P}(S_t = j \mid S_{t-1} = i)$ is the probability of transitioning from regime $i$ to regime $j$. The transition probabilities are collected in the $K \times K$ row-stochastic regime transition matrix
    \begin{equation}
        P=
        \begin{bmatrix}
            p_{11} & p_{12} & \cdots & p_{1K} \\
            p_{21} & p_{22} & \cdots & p_{2K} \\
            \vdots & \vdots & \ddots & \vdots \\
            p_{K1} & p_{K2} & \cdots & p_{KK}
        \end{bmatrix},
        \qquad \sum_{j=1}^{K} p_{ij} = 1 \quad \text{for all } i.
    \end{equation}
\end{definition}

The observation distribution admits a density function, which we denote by $f(y_t \mid x_t, \theta) = f(y_t \mid h_t)$, as it does not depend on the regime $s_t$ or on any of the model parameters $\theta$. 
Note that mathematically speaking the latent state $H_t$ is the log variance of the return $Y_t$, but it is often referred to as the log volatility in the literature \cite{carvalho2007simulation, dahlin2019getting, bao2018particle}. We will also follow this convention and refer to $H_t$ as the log volatility.

Observe that, given the regimes, the log volatility process is essentially an AR(1) process with regime-specific means. The initial distribution of $H_0$ is chosen as the stationary distribution of the corresponding AR(1) component, and the initial regime has a uniform prior \cite{so1998stochastic}. The initial distribution admits a density function, which can be written as $\chi(x_0 \mid \theta) = \chi(s_0 \mid \theta) \chi(h_0 \mid s_0, \theta)$.

At first, it may seem that such a conditional definition breaks the Markov property of the latent state process, as $H_t$ depends on $S_t$. The following proposition shows that this is not the case.

\begin{proposition}
    \label{prop:mssv_markov}
    The latent process $\{X_t\}_{t \geq 0} = \{(H_t, S_t)\}_{t \geq 0}$ of the MSSV model is a Markov process with transition density
    \begin{equation}
        p(x_t \mid x_{t-1}, \theta) = p(s_t \mid s_{t-1}, \theta)\, p(h_t \mid h_{t-1}, s_t, \theta).
    \end{equation}
\end{proposition}
\begin{proof}
    Using the definition of conditional probability
    \begin{align}
        p(x_t \mid x_{t-1}, \theta) &= p((h_t, s_t) \mid (h_{t-1}, s_{t-1}), \theta) \\
        &= p(h_t \mid h_{t-1}, s_{t-1}, s_t, \theta)\, p(s_t \mid h_{t-1}, s_{t-1}, \theta) \\
        &= p(h_t \mid h_{t-1}, s_t, \theta)\, p(s_t \mid s_{t-1}, \theta).
    \end{align}
\end{proof}

In order for the MSSV model to be well-defined, the following assumptions are made.
\begin{enumerate}
    \item The persistence $\phi$ lies in the range $(-1, 1)$.
    \item Each regime corresponds to at least one time point.
    \item The regime means are ordered such that $\mu_1 < \mu_2 < \ldots < \mu_K$.
\end{enumerate}

Assumption 1 ensures that the AR(1) process is stationary, as shown in the next proposition. This is needed to ensure that the log volatility does not diverge to infinity. Assumption 2 is necessary to avoid empty regimes, as otherwise some $\mu_i$ would be unidentified.
Assumption 3 is a common identification constraint in mixture models, which prevents label switching and ensures that the regimes are uniquely defined \cite{so1998stochastic}. 

In this thesis, we ensure the identifiability of the regimes using the following parameterization:
\begin{equation}
    \mu_i = \mu_{i-1} + e^{\delta_i}, \qquad i = 2, \ldots, K,
\end{equation}
where $\delta_i \in \mathbb{R}$. This parameterization is motivated by the fact that the exponential function is strictly positive, which guarantees that $\mu_i > \mu_{i-1}$ while leaving $\delta_i$ unconstrained. 

The parameters of the MSSV model are then given by the vector
\begin{equation}
    \theta = \{\mu_1, \delta_2, \ldots, \delta_K, \phi, \eta^2, P\}.
\end{equation}
Since $P$ is row stochastic, it has $K(K-1)$ free parameters. Therefore, the total number of free parameters is $K + 2 + K(K-1) = K^2 + 2$.

\begin{remark}
    In the case of $K=1$, the MSSV model reduces to the standard stochastic volatility (SV) model, in which only the log volatility $H_t$ is a latent variable and there are no regimes.
\end{remark}

\section{Problem formulation}
\label{sec:problem_formulation}
The goal of this thesis is to estimate volatility in financial markets using the MSSV model with parameters $\theta$. 
Let $y_{1:t} = \{y_1, y_2, \ldots, y_t\}$ denote the sequence of observed returns up to time $t \leq T$. Similarly, let $x_{0:t} = \{x_0, x_1, \ldots, x_t\}$ denote the sequence of latent states up to time $t$, where $x_0$ is the initial prior state.
In general, we can distinguish between two versions of the problem.
\begin{enumerate}
    \item \textbf{With known parameters.} If we assume that the model parameters $\theta$ are known, then the main objective is the smoothing distribution $p(x_{0:t} \mid y_{1:t}, \theta)$.
    \item \textbf{With unknown parameters.} In reality, the parameters $\theta$ of the MSSV model are unknown, so we would also like to estimate them. The main objective is then the joint posterior distribution $p(x_{0:t}, \theta \mid y_{1:t})$, from which we can compute the smoothing, filtering and predictive distributions of the latent states and the posterior distribution of the parameters, as in Definition~\ref{def:distributions_of_interest}.
\end{enumerate}
Both targets are determined, up to normalizing constants, by the joint density of the latent states and observations.

\begin{proposition}
    \label{prop:joint_density}
    For the MSSV model, the joint density of the latent states and observations is
    \begin{equation}
        \label{eq:joint_density}
        p(x_{0:t}, y_{1:t} \mid \theta) = \chi(x_0 \mid \theta) \prod_{i=1}^{t} p(x_i \mid x_{i-1}, \theta) f(y_i \mid x_i).
    \end{equation}
    Consequently, with known parameters,
    \begin{equation}
        \label{eq:joint_posterior_known_parameters}
        p(x_{0:t} \mid y_{1:t}, \theta) \propto \chi(x_0 \mid \theta) \prod_{i=1}^{t} p(x_i \mid x_{i-1}, \theta) f(y_i \mid x_i),
    \end{equation}
    and, given a prior $p(\theta)$, with unknown parameters,
    \begin{equation}
        \label{eq:joint_posterior_unknown_parameters}
        p(x_{0:t}, \theta \mid y_{1:t}) \propto p(\theta) \chi(x_0 \mid \theta) \prod_{i=1}^{t} p(x_i \mid x_{i-1}, \theta) f(y_i \mid x_i).
    \end{equation}
\end{proposition}
\begin{proof}
    By the definition of conditional density,
    \begin{equation}
        p(x_{0:t}, y_{1:t} \mid \theta) = p(y_{1:t} \mid x_{0:t}, \theta)\, p(x_{0:t} \mid \theta).
    \end{equation}
    The first factor can be decomposed as
    \begin{align}
        p(y_{1:t} \mid x_{0:t}, \theta) &= p(y_t \mid y_{1:t-1}, x_{0:t}, \theta) p(y_{1:t-1} \mid x_{0:t}, \theta) \\
        &= f(y_t \mid x_t, \theta) p(y_{1:t-1} \mid x_{0:t-1}, \theta) \\
        &= \prod_{i=1}^{t} f(y_i \mid x_i),
    \end{align}
    where we used the fact that each observation $y_i$ is conditionally independent of the other observations and states given $x_i$ (Definition~\ref{def:state_space_model}).

    Similarly, the second factor can be decomposed as
    \begin{equation}
        p(x_{0:t} \mid \theta) = \chi(x_0 \mid \theta) \prod_{i=1}^{t} p(x_i \mid x_{0:i-1}, \theta) = \chi(x_0 \mid \theta) \prod_{i=1}^{t} p(x_i \mid x_{i-1}, \theta).
    \end{equation}
    Multiplying both factors yields~\eqref{eq:joint_density}.

    For known parameters, Bayes' theorem gives
    \begin{equation}
        p(x_{0:t} \mid y_{1:t}, \theta) = \frac{p(x_{0:t}, y_{1:t} \mid \theta)}{p(y_{1:t} \mid \theta)},
    \end{equation}
    where the denominator does not depend on $x_{0:t}$, which yields~\eqref{eq:joint_posterior_known_parameters}. For unknown parameters,
    \begin{equation}
        p(x_{0:t}, \theta \mid y_{1:t}) = \frac{p(x_{0:t}, y_{1:t} \mid \theta)\, p(\theta)}{p(y_{1:t})},
    \end{equation}
    where the denominator does not depend on $(x_{0:t}, \theta)$, which yields~\eqref{eq:joint_posterior_unknown_parameters}.
\end{proof}

Equations~\eqref{eq:joint_posterior_known_parameters} and~\eqref{eq:joint_posterior_unknown_parameters} both reveal the sequential nature of the problem. If we initially sample the prior latent state $x_0$ (and the model parameters $\theta$, if they are unknown), we can sequentially propagate the latent states through the transition density and update them with new observations through the observation density. This is the main idea behind sequential Monte Carlo methods, which are presented in the next chapter.

\subsection{Prior distributions of the parameters}
\label{sec:priors}
In the case of unknown model parameters $\theta$, we need to specify their prior distributions. 
We assume that the priors of the parameters are independent, so that the joint prior distribution can be factorized as
\begin{equation}
    p(\theta) = p(\mu_1) \left(\prod_{i=2}^{K} p(\delta_i)\right) p(\phi)\, p(\eta^2)\, p(P).
\end{equation}

Even though the model is defined in terms of the parameters $\delta_i$ and $\phi$, we define the priors in terms of $\Delta_i = \mu_i - \mu_{i-1} = e^{\delta_i}$ and $\varphi = \frac{\phi + 1}{2}$. 
The reason for reparameterizing $\phi$ as $\varphi$ is that $\varphi$ lies in the range $(0, 1)$, which allows us to use the Beta distribution as a prior. 
The reason for reparameterizing $\delta_i$ as $\Delta_i$ is that we can then use a truncated normal distribution as a prior, which is more convenient for conditional sampling later on.  

We choose the following prior distributions for the model parameters:
\begin{align}
    \mu_1 &\sim \mathcal{N}(a_0,5^2),\\
    \Delta_i &\sim \mathcal{TN}_{(0,\infty)}(0,5^2), \qquad i=2,\ldots,K,\\
    \varphi &\sim \mathcal{B}(6,2),\\
    \eta^2 &\sim \mathcal{IG}(2.001,0.1),\\
    P_{i,:} &\sim \mathcal{D}(\alpha_{i,1},\ldots,\alpha_{i,K}), \qquad i=1,\ldots,K,
\end{align}
where $\alpha_{i,j}=1.5$ for $i\neq j$ and $\alpha_{i,i}=40$. Here, $\mathcal{TN}_{(0,\infty)}$ denotes the truncated normal distribution with support $(0,\infty)$, $\mathcal{B}$ denotes the Beta distribution, $\mathcal{IG}$ denotes the inverse gamma distribution, and $\mathcal{D}$ denotes the Dirichlet distribution.

The priors are chosen to be informative while allowing sufficient flexibility. For $\mu_1$, $a_0$ is determined through prior elicitation, while the variance $5^2$ allows for substantial uncertainty. The prior for $\Delta_i$ allows for a wide range of possible values, but still encourages the regime means not to be too far apart. The $\mathcal{B}(6,2)$ prior for $\varphi$ favors positive persistence, as is common in financial econometrics. Additionally, previous results from \cite{pan2013testing} suggest that the choice of prior for this parameter has little effect on the posterior, so we chose hyperparameters that lie in between the values considered there.
For $\eta^2$, we use $\mathcal{IG}(2.001,0.1)$, following a combination of the values considered in \cite{carvalho2007simulation,pan2013testing}. This prior has a large variance, allowing for a wide range of values.
For the transition matrix $P$, we depart from the commonly used symmetric Dirichlet prior with hyperparameters $0.5$ \cite{so1998stochastic,carvalho2007simulation,pan2013testing}. Such a prior places substantial mass near $0$ and $1$, which may encourage regimes to either disappear or become excessively persistent. Instead, motivated by \cite{shibata2005bayesian}, we use $\alpha_{i,i}=40$ and $\alpha_{i,j}=1.5$ for $i\neq j$. In the two-regime case, this corresponds to a $\mathcal{B}(40,1.5)$ prior for the diagonal transition probabilities, with mean $0.964$, implying an expected regime duration of approximately 28 days. We find this expected duration reasonable for financial data.
